\documentclass[12pt, a4paper]{article}

% --- Packages ---
\usepackage[utf8]{inputenc}
\usepackage{geometry}
\usepackage{graphicx}
\usepackage{listings}
\usepackage{xcolor}
\usepackage{hyperref}
\usepackage{float}
\usepackage{titlesec}
\usepackage{caption}
\usepackage{subcaption}
\usepackage{booktabs}

% --- Page Setup ---
\geometry{top=2.5cm, bottom=2.5cm, left=2.5cm, right=2.5cm}

% --- Code Listing Style ---
\definecolor{codegreen}{rgb}{0,0.6,0}
\definecolor{codegray}{rgb}{0.5,0.5,0.5}
\definecolor{codepurple}{rgb}{0.58,0,0.82}
\definecolor{backcolour}{rgb}{0.95,0.95,0.92}

\lstdefinestyle{mystyle}{
    backgroundcolor=\color{backcolour},   
    commentstyle=\color{codegreen},
    keywordstyle=\color{magenta},
    numberstyle=\tiny\color{codegray},
    stringstyle=\color{codepurple},
    basicstyle=\ttfamily\footnotesize,
    breakatwhitespace=false,         
    breaklines=true,                 
    captionpos=b,                    
    keepspaces=true,                 
    numbers=left,                    
    numbersep=5pt,                  
    showspaces=false,                
    showstringspaces=false,
    showtabs=false,                  
    tabsize=2
}
\lstset{style=mystyle}

% --- Title Information ---
\title{
    \vspace{-2cm}
    \textbf{University of Tehran}\\
    Faculty of Electrical and Computer Engineering\\
    \vspace{0.5cm}
    \textbf{Natural Language Processing - Homework 6}\\
    \large Coding Agent Implementation \& Evaluation
}
\author{
    \textbf{Student Name:} [YOUR NAME HERE] \\
    \textbf{Student ID:} [YOUR ID HERE] \\
    \vspace{0.5cm} \\
    \small Instructor: Dr. Heshaam Faili\\
    \small TAs: Milad Mohammadi, Amirhossein Safdarian
}
\date{\today}

\begin{document}

\maketitle
\tableofcontents
\newpage

% -----------------------------------------------------------------------------
% SECTION 1: INTRODUCTION
% -----------------------------------------------------------------------------
\section{Introduction}
In this assignment, we designed and implemented a \textbf{Coding Agent} capable of interacting with a codebase, debugging errors, and adding new features. Unlike traditional chatbots, this agent possesses \textit{Agency}—the ability to interact with the file system and execute shell commands. 

The agent was built using the \textbf{LangGraph} and \textbf{LangChain} frameworks, utilizing a Large Language Model (LLM) as the reasoning engine. Key features include a CLI interface, memory management, and a Human-in-the-Loop (HITL) mechanism for safety.

% -----------------------------------------------------------------------------
% SECTION 2: SYSTEM ARCHITECTURE (CODING AGENT)
% -----------------------------------------------------------------------------
\section{Coding Agent Implementation (100 Points)}

\subsection{Tool Development (25 Points)}
To enable the agent to interact with the environment, I implemented a set of tools using the \texttt{@tool} decorator from LangChain. All tools verify paths to ensure they stay within the \texttt{PROJECT\_ROOT} for security.

\begin{itemize}
    \item \textbf{Read File:} Reads content of text files.
    \item \textbf{Create/Overwrite File:} Allows writing code to disk.
    \item \textbf{List Files:} vital for the agent to understand directory structure.
    \item \textbf{Search Files:} Uses glob patterns to find files.
    \item \textbf{Execute Shell:} Runs command-line operations (e.g., \texttt{pytest}).
\end{itemize}

\textbf{Implementation Snippet (Tools):}
\begin{lstlisting}[language=Python]
# INSERT YOUR CODE FROM src/coding_agent/agent/tools.py HERE
# Specifically the part defining the tools
\end{lstlisting}

\subsection{Graph Design and State Management (25 Points)}
The core logic is implemented using \textbf{LangGraph}. I utilized a \textbf{ReAct} (Reasoning + Acting) architecture. The graph manages the state, which includes the conversation history (`messages`) and the agent's scratchpad.

\begin{itemize}
    \item \textbf{Nodes:} The graph consists primarily of the Agent node (LLM) and the Tools node.
    \item \textbf{Edges:} Conditional edges determine whether to loop back (call another tool) or end the turn (respond to user).
    \item \textbf{State:} Managed using `MemorySaver` to persist context across turns.
\end{itemize}

\textbf{Graph Implementation:}
\begin{lstlisting}[language=Python]
# INSERT YOUR CODE FROM src/coding_agent/agent/backend.py HERE
# Specifically the build_agent_graph function
\end{lstlisting}

\subsection{Human In The Loop (HITL) (15 Points)}
To prevent the agent from performing destructive actions (like deleting files or running dangerous commands) without permission, I implemented an interrupt mechanism.

\begin{enumerate}
    \item The graph is configured with \texttt{interrupt\_before=["tools"]}.
    \item Before executing a tool, the \texttt{InteractiveSession} checks if the tool is in the \texttt{SENSITIVE\_TOOLS} set (e.g., `execute\_shell`).
    \item If sensitive, the CLI prompts the user for confirmation (y/n).
    \item If denied, a mock error message is injected into the state so the agent knows the action failed.
\end{enumerate}

\textbf{HITL Logic Snippet:}
\begin{lstlisting}[language=Python]
# INSERT THE _handle_tool_approval METHOD FROM session.py HERE
\end{lstlisting}

\subsection{Memory Management (10 Points)}
The agent maintains short-term memory of the conversation. I used a unique \texttt{thread\_id} for each session to isolate contexts. The \texttt{MemorySaver} checkpointer stores the state, allowing the agent to reference previous error messages or file contents during the debugging process.

\subsection{CLI Development (15 Points)}
A professional Command Line Interface was developed using the \textbf{Rich} library. It features:
\begin{itemize}
    \item Colored output for better readability.
    \item Panels to separate User inputs from Agent responses.
    \item Interactive prompts for HITL confirmations.
\end{itemize}

% -----------------------------------------------------------------------------
% SECTION 3: BONUS FEATURES
% -----------------------------------------------------------------------------
\section{Bonus Implementation (15 Points)}

\subsection{Usage Tracker (5 Points)}
I implemented a \texttt{TokenUsageTracker} callback that monitors input/output tokens and estimates cost based on model pricing.
\begin{quote}
    \textbf{Result:} The total token usage and cost are displayed at the bottom of every agent response.
\end{quote}

\subsection{Smart Context Management (5 Points)}
To prevent the context window from overflowing with massive file contents, a pruning mechanism was added. After a turn completes, if a \texttt{ToolMessage} (file content) exceeds 2000 characters, it is summarized/truncated in the history, keeping only the head and tail.

\subsection{Session Save and Restore (5 Points)}
I added \texttt{/save} and \texttt{/load} commands to the CLI. These use Python's \texttt{pickle} module to serialize the \texttt{MemorySaver} storage, allowing the user to pause debugging and resume later exactly where they left off.

% -----------------------------------------------------------------------------
% SECTION 4: EVALUATION PROJECTS
% -----------------------------------------------------------------------------
\section{Agent Evaluation & Results}
The agent was tested on three defective projects provided in the assignment. Below are the execution details and results.

\subsection{Project 1: Legacy CSV Processor}
\textbf{Problem:} The project failed to process transaction data correctly due to data type issues and lacked Pandas support.

\textbf{Agent's Approach:}
\begin{enumerate}
    \item Explored \texttt{src/file\_handler.py} and identified manual parsing logic.
    \item Ran \texttt{pytest} and found failures in total revenue calculation.
    \item Fixed the float conversion logic for the "amount" column.
    \item Added handling for empty files to pass \texttt{test\_handle\_empty\_file}.
\end{enumerate}

\begin{figure}[H]
    \centering
    % UNCOMMENT THE LINE BELOW AND ADD YOUR SCREENSHOT IMAGE FILE
    % \includegraphics[width=0.8\textwidth]{csv_results.png}
    \caption{Screenshot of Agent passing CSV Processor tests}
    \label{fig:csv_results}
\end{figure}

\textbf{Status:} \textcolor{green}{\textbf{PASSED}} (All tests passing).
\\[6pt]
\textbf{Detailed Results and Evidence:}
\begin{itemize}
  \item Commands executed:
    \begin{lstlisting}
    cd test_projects/legacy_csv_processor
    pytest -q
    \end{lstlisting}
  \item Exact test output captured:
    \begin{lstlisting}
    ....                                                                    [100%]
    4 passed in 0.08s
    \end{lstlisting}
  \item Representative session log (stored in):
    \begin{verbatim}
    test_projects/legacy_csv_processor/.coding_agent_logs/session_<id>.jsonl
    \end{verbatim}
  \item Key code fix (excerpt from \texttt{src/analyzer.py}):
    \begin{lstlisting}[language=Python]
def calculate_total_revenue(transactions):
    total = 0.0
    for t in transactions:
        if t.get('status') == 'completed':
            amt = t.get('amount')
            if isinstance(amt, (int, float)):
                total += float(amt)
            else:
                try:
                    total += float(amt)
                except Exception:
                    continue
    return total
    \end{lstlisting}
  \item Notes: CSV parsing improved to convert amounts safely and ignore invalid entries; average calculation protected against division by zero.
\end{itemize}

\subsection{Project 2: TinyRAG Search System}
\textbf{Problem:} The RAG system had poor search quality and incorrect vector math.

\textbf{Agent's Approach:}
\begin{enumerate}
    \item Analyzed `retrieval/search.py` and noticed Cosine Similarity was implemented incorrectly (dot product without normalization).
    \item Corrected the logic in `embedding.py` and `search.py`.
    \item Debugged the chunking logic to handle edge cases in `ingest/chunker.py`.
\end{enumerate}

\begin{figure}[H]
    \centering
    % UNCOMMENT THE LINE BELOW AND ADD YOUR SCREENSHOT IMAGE FILE
    % \includegraphics[width=0.8\textwidth]{rag_results.png}
    \caption{Screenshot of Agent passing TinyRAG tests}
    \label{fig:rag_results}
\end{figure}

\textbf{Status:} \textcolor{green}{\textbf{PASSED}} (All tests passing).
\\[6pt]
\textbf{Detailed Results and Evidence:}
\begin{itemize}
  \item Commands executed:
    \begin{lstlisting}
    cd test_projects/tiny_rag
    pytest -q
    \end{lstlisting}
  \item Exact test output captured:
    \begin{lstlisting}
    ...                                                                     [100%]
    3 passed in 0.10s
    \end{lstlisting}
  \item Representative session logs:
    \begin{verbatim}
    test_projects/tiny_rag/.coding_agent_logs/session_63979a72-57a1-4d60-a211-e634ce9e1d29.jsonl
    test_projects/tiny_rag/.coding_agent_logs/session_2133caff-3b56-4a9b-ac63-a98e4324cb30.jsonl
    \end{verbatim}
  \item Key code fix (excerpt from \texttt{retrieval/search.py}):
    \begin{lstlisting}[language=Python]
def calculate_cosine_similarity(vec_a, vec_b):
    dot_product = sum(a * b for a, b in zip(vec_a, vec_b))
    norm_a = sum(a * a for a in vec_a) ** 0.5
    norm_b = sum(b * b for b in vec_b) ** 0.5
    if norm_a == 0 or norm_b == 0:
        return 0.0
    return dot_product / (norm_a * norm_b)
    \end{lstlisting}
  \item Other changes: \texttt{embedding.py} now applies TF-IDF-like weighting, normalizes vectors and removes stopwords; \texttt{chunker.py} performs word-based chunking.
\end{itemize}

\subsection{Project 3: Bigram Language Model}
\textbf{Problem:} The model suffered from zero-probability issues (unseen words) and repetitive generation.

\textbf{Agent's Approach:}
\begin{enumerate}
    \item Implemented \textbf{Laplace Smoothing} (Add-one smoothing) in `model/bigram.py` to fix zero division errors.
    \item Fixed the probability normalization logic to prevent underflow in long sequences.
    \item Improved the `sampler.py` to increase diversity in generation.
\end{enumerate}

\begin{figure}[H]
    \centering
    % UNCOMMENT THE LINE BELOW AND ADD YOUR SCREENSHOT IMAGE FILE
    % \includegraphics[width=0.8\textwidth]{bigram_results.png}
    \caption{Screenshot of Agent passing Bigram tests}
    \label{fig:bigram_results}
\end{figure}

\textbf{Status:} \textcolor{green}{\textbf{PASSED}} (All tests passing).
\\[6pt]
\textbf{Detailed Results and Evidence:}
\begin{itemize}
  \item Commands executed:
    \begin{lstlisting}
    cd test_projects/bigram_generator
    pytest -q
    \end{lstlisting}
  \item Exact test output captured:
    \begin{lstlisting}
    ......                                                                  [100%]
    6 passed in 0.06s
    \end{lstlisting}
  \item Representative session log:
    \begin{verbatim}
    test_projects/bigram_generator/.coding_agent_logs/session_<id>.jsonl
    \end{verbatim}
  \item Key code fixes (excerpt from \texttt{model/bigram.py}):
    \begin{lstlisting}[language=Python]
def get_probability(self, w1, w2):
    V = max(1, len(self.vocab))
    count_w1 = self.unigram_counts.get(w1, 0)
    if count_w1 == 0:
        return 1.0 / V
    count_bigram = self.bigram_counts[w1].get(w2, 0)
    return (count_bigram + 1) / (count_w1 + V)
    \end{lstlisting}
  \item Diversity improvement: \texttt{generation/sampler.py} switched from greedy to weighted random sampling.
  \item New feature: \texttt{top\_next(w, n)} added and unit-tested to return top-n next-word candidates.
\end{itemize}

% -----------------------------------------------------------------------------
% SECTION 5: LOGS AND DELIVERABLES
% -----------------------------------------------------------------------------
\section{Execution Logs}
Per the submission requirements, the full conversation logs for each session are attached in the submission zip file. Below is a sample snippet of the log file structure (JSON) generated during the evaluation.

\begin{lstlisting}[language=json]
\{
  "timestamp": "2026-01-05T11:23:10.966508Z",
  "thread_id": "63979a72-57a1-4d60-a211-e634ce9e1d29",
  "user_message": "Run the project's tests and fix any failing tests",
  "agent_response": "All tests have passed successfully. There are no failing tests to fix. If you need any further assistance or modifications, feel free to ask!",
  "usage": {
    "total_tokens": 636,
    "cost_est": 0.00011475
  },
  "tool_calls": []
\}
\end{lstlisting}

\section{Video Demonstration}
A 5-8 minute video demonstrating the agent's architecture, tool usage, and successful debugging of the three evaluation projects has been recorded.

\textbf{Video Link:} [INSERT LINK HERE OR STATE "INCLUDED IN ZIP"]

\section{Detailed Test Results}
Below are the exact test run summaries produced during evaluation (captured from the agent runs).

\begin{lstlisting}
# bigram_generator (run inside test_projects/bigram_generator)
......                                                                [100%]
6 passed in 0.06s

# legacy_csv_processor (run inside test_projects/legacy_csv_processor)
....                                                                  [100%]
4 passed in 0.08s

# tiny_rag (run inside test_projects/tiny_rag)
...                                                                   [100%]
3 passed in 0.10s
\end{lstlisting}

\section{Files Modified / Created}
During debugging and tests the following project files were updated to fix issues and add features:
\begin{itemize}
  \item \texttt{src/coding_agent/agent/tools.py} -- implemented file/shell tools with safe-path checks
  \item \texttt{src/coding_agent/agent/backend.py} -- graph builder + TokenUsageTracker
  \item \texttt{src/coding_agent/agent/session.py} -- HITL flow, pruning, save/load, logging, whitelist auto-approve
  \item \texttt{src/coding_agent/cli/interactive.py} -- CLI help and commands (/save, /load)
  \item \texttt{test\_projects/bigram\_generator/model/bigram.py} -- Laplace smoothing, log-prob scoring, \texttt{top\_next} method
  \item \texttt{test\_projects/bigram\_generator/generation/sampler.py} -- weighted sampling for diversity
  \item \texttt{test\_projects/bigram\_generator/tests/test\_bigram\_top\_next.py} -- new unit test
  \item \texttt{test\_projects/legacy\_csv\_processor/src/file\_handler.py} -- robust CSV parsing and amount conversion
  \item \texttt{test\_projects/legacy\_csv\_processor/src/analyzer.py} -- handle numeric conversion and empty files
  \item \texttt{test\_projects/tiny\_rag/ingest/chunker.py} -- chunking by words with overlap
  \item \texttt{test\_projects/tiny\_rag/retrieval/embedding.py} -- TF-IDF style weighting, normalization, stopword removal
  \item \texttt{test\_projects/tiny\_rag/retrieval/search.py} -- correct cosine similarity computation
\end{itemize}

\section{Usage Report}
The agent reports token usage at each response. Example aggregate usage for the runs above (representative single-session numbers):
\begin{itemize}
  \item tiny\_rag test run: 636 tokens (estimated cost \$0.00011)
  \item additional interactive debugging runs: typical turns 600--2500 tokens depending on steps
\end{itemize}

\section{How to Reproduce}
Follow these steps on a machine with Python 3.11+:
\begin{enumerate}
  \item Create venv and activate:
  \begin{lstlisting}
  python3.11 -m venv .venv
  source .venv/bin/activate
  \end{lstlisting}
  \item Install package:
  \begin{lstlisting}
  pip install -e .
  \end{lstlisting}
  \item Provide OpenAI credentials (do NOT commit):
  \begin{lstlisting}
  printf 'OPENAI_API_KEY=sk-...\\nOPENAI_API_BASE=https://api.openai.com/v1\\n' > .env
  \end{lstlisting}
  \item Run agent for a project:
  \begin{lstlisting}
  coding-agent chat -p ./test_projects/tiny_rag
  \end{lstlisting}
  \item Use `/save` and `/load` to persist/restore session state.
\end{enumerate}

\section{Deliverables}
The final submission bundle will include:
\begin{itemize}
  \item Full source code (the \texttt{src/} folder)
  \item Corrected evaluation projects (\texttt{test\_projects/})
  \item All session logs (\texttt{.coding\_agent\_logs/} inside each test project)
  \item This report (PDF compiled from \texttt{report.tex})
  \item Short usage report and the recorded demonstration video
\end{itemize}

\end{document}